% ch13: 原文 GLDraft_D.tex 行 3170–3425；§9 + 致谢
% 只译散文（标题、正文、图题、脚注）；公式、\label、\ref、\eqref、\cite、图文件名保持不变。

\section{来自辛玻色子的波函数} \label{sec:symp}

\subsection{辛玻色子的基础}\label{sb}
如我们在第 \ref{Enriched} 节所讨论的，由三维超多重态增广的 Neumann 边界条件所产生的态，可由一个路径积分描述
\begin{equation}
\Psi[a] = \int DZ D\bar Z\, e^{\int_C \left[\langle Z, \bar \partial_a Z \rangle -  \langle 
\bar Z, \partial_a \bar Z \rangle\right]}
\end{equation} 
这是{\it 辛玻色子}（symplectic bosons）路径积分的非手征版本。手征辛玻色子是手征费米子的 Grassmann 偶类比物。它们是 $\beta \gamma$ 系统的一个特例，其中 $\beta$ 与 $\gamma$ 的共形维数都被取为 $1/2$。\zy{译者按：$\beta\gamma$ 系统是共形维数 $(\lambda,1-\lambda)$ 的自由玻色鬼场；取 $\lambda=1/2$ 时两个场都是自旋 $1/2$ 的玻色场，其 OPE 因辛形式反对称而与费米子相似，统计却是玻色。标准处理见 J. Polchinski, \emph{String Theory}, Vol.~1 (1998) 第 2 章；D. Friedan, E. Martinec, S. Shenker, \emph{Nucl.\ Phys.} B271 (1986) 93。}

具体地说，手征辛玻色子是一组 $2 n$ 个二维自旋 $1/2$ 的手征玻色场 $Z^a$，其作用量为 
\begin{equation}
\int_{C}\left( \omega_{ab} Z^a \bar\partial Z^b + {\cal A}_{ab} Z^a Z^b\right)
\end{equation}
其中 $\omega_{ab}$ 是常值辛形式，我们还加入了与一个 $(0,1)$ 型背景联络 ${\cal A}_{ab}$ 的耦合，该联络在黎曼曲面 $C$ 上定义一个 $\Sp(2n)$ 丛。\footnote{由于辛玻色子是旋量，我们严格说来并不需要分别定义一个旋量结构与一个 $\Sp(2n)$ 丛。我们改为可以指定一个 $\Spin\cdot \Sp(2n)$ 丛，这一概念与第 \ref{toposubt} 节的 $\Spin\cdot \SU(2)$ 丛完全类似。} 本节除非另有说明，我们采用 Einstein 求和约定。 
 
 与手征费米子的类比并不完全。手征费米子是定义良好的二维 (spin)CFT。
手征辛玻色子则有轻微的反常。反常表现为手征配分函数的一个符号不确定性
\begin{equation}
\int DZ \,e^{\int_{C}\left( \omega_{ab} Z^a \bar\partial Z^b + {\cal A}_{ab} Z^a Z^b\right)} = \frac{1}{\sqrt{\det \bar \partial_{\cal A} }}
\end{equation}
其中我们把作用在 $K_C^{1/2} \otimes E$ 的截面上的 $\bar \partial$ 算子记作 $\bar \partial_{\cal A}$，而 $E$ 是附属于该 $\Sp(2n)$ 丛的秩 $2n$ 丛。\footnote{注意，一般而言 $\bar \partial_{\cal A}$ 没有零模，泛函行列式是定义良好的。对于出现零模的特殊 $\Sp(2n)$ 丛选取，配分函数发散。} 
不过，非手征配分函数
\begin{equation}
\int DZ D\bar Z\, e^{\int_C \left[\langle Z, \bar \partial_A Z \rangle -  \langle 
\bar Z, \partial_A \bar Z \rangle\right]} =  \frac{1}{|\det \bar \partial_{\cal A} |}
\end{equation}
是无歧义的：它可以沿循环 $\bar Z{}^a = (Z^a)^*$ 作真正的路径积分来定义（也就是说，该积分循环由如下要求给出：$\bar Z{}^a$ 是 $Z^a$ 的复共轭）。\zy{译者按（无外部依据）：原文此句作 “by an actual the path integral”，多一冠词 the；§9.2 又有 “the the kinetic operator” 重复 the；§9.4 开头列举三个谱流算符时，$\Sigma'$ 与 $\Sigma''$ 的 $\bar\mu$ 缺撇（其下分项中则各自带撇）。以上均为原文笔误，译文照原文处理，未径改。}

我们在此记下算子乘积展开（OPE）
\begin{equation}
Z^a(z) Z^b(w) \sim \frac{\omega^{ab}}{z-w},
\end{equation}
其中 $\omega^{ab}$ 是辛形式的逆。相应的模代数为 
\begin{equation}
[Z^a_m, Z^b_n] = \omega^{ab} \delta_{n+m,0}.
\end{equation}
由于自旋为半整数，在 Neveu–Schwarz 扇区中——特别是在真空模中——模指标 $n$、$m$ 为半整数。真空满足 
\begin{equation}
Z^a_n |0\rangle =0 \qquad \qquad n>0.
\end{equation}
而模指标在 Ramond 扇区模中为整数，我们将在第 \ref{sec:Ramond} 节讨论这些模。

作用量对 $\Sp(2n)$ 联络 ${\cal A}$ 的变分给出 level 为 $- \frac12$ 的 Kac–Moody 流
\begin{equation}
J^{ab} = \frac12 :Z^a Z^b:
\end{equation}
分数 level 是手征理论整体反常的又一表现。\footnote{$J^{ab}$ 流实际上生成 $\widehat{\mathfrak{sp}}(2n)_{-\frac12}$ 手征代数的一个商：流的双线性型的某些线性组合以及流的导数会消失。辛玻色子真空模中 level $2$ 后代的数目，小于 Kac–Moody 真空模中 level $2$ 后代的数目。} 非手征理论的配分函数定义了 $\Sp(2n,\C)$ 丛空间上一个丛 $|{\cal L}|$ 的定义良好的截面，其中 ${\cal L}$ 是对应于 Kac–Moody level $-1$ 的线丛。 

一旦我们特殊化到四维理论的规范群 $G \subset \Sp(2n)$，就会得到 level 未必临界的 $G$ 流。这标志着一种规范反常，它阻碍变形 Neumann 边界条件与增广 Neumann 边界条件之间存在二维交汇。若 level 比临界 level 更负，我们可以尝试用某个辅助二维系统（例如一组自由费米子）来抵消该反常。在下文即将讨论的例子中，反常从一开始就不存在。

\subsection{三基本例子}\label{tri}

现在我们特殊化到 $n=4$，并集中于 $\Sp(8)$ 的一个子群 $G\equiv \SL(2) \times \SL(2)\times \SL(2)$。
换言之，我们把 $\mathbb{C}^8$ 与 $\mathbb{C}^2 \otimes\mathbb{C}^2 \otimes\mathbb{C}^2$ 等同起来，
并把理论耦合到 $\SL(2) \times \SL(2)\times \SL(2)$ 的联络 $a$。

这给出一些增广 Neumann 边界条件，人们猜测它们与一个三对角界面（tri-diagonal interface）$S$-对偶，后者是支撑在 $\M_H \times \M_H \times \M_H$ 的对角线上、带有平凡 $\CP$ 丛的 BBB 膜 \cite{Gaiotto:2016hvd,Benini:2010uu}。\footnote{这一 $S$-对偶陈述可以推广到其他 $G$，或推广到三个以上副本之间的对角界面。相应的增广 Neumann 边界条件使用 \cite{Benini:2010uu} 中定义的理论。这些理论的边界手征代数已由 Arakawa 的工作 \cite{Arakawa:2018egx} 得知，它们具有临界 level 的 $G$ 流。相应的态将是由 Arakawa 的手征代数构造的非手征二维共形场论的配分函数，而后者目前尚不清楚。} 这一 $S$-对偶等同的基本推论是：由配分函数产生的态 $\Psi[a]$ 应当把三份 $\SL(2)$ Hitchin 哈密顿量与量子 Hecke 算符的作用交织起来。本节余下部分的目标就是证实这一点。
 
我们把辛玻色子场记作 $Z^{\alpha \beta \gamma}(z)$，把辛形式记作
$\epsilon_{\alpha \alpha'}\epsilon_{\beta \beta'}\epsilon_{\gamma \gamma'}$。我们得到三份
$\widehat{\mathfrak{sl}}(2)_{-2}$ Kac–Moody 流，例如
\begin{equation}
J^{\alpha \alpha'} =\frac12 \epsilon_{\beta \beta'}\epsilon_{\gamma \gamma'} :Z^{\alpha \beta \gamma}Z^{\alpha' \beta' \gamma'}:
\end{equation}
我们把这三组流简单地记作 $J$、$J'$、$J''$，在可能时省略指标。 

关键在于，这些流具有临界 level。相应地，Sugawara 向量是中心元。一个值得注意的观察是：此处三个 Sugawara 向量实际上彼此重合：
\begin{equation} \label{eq:sugathree}
:JJ: = :J'J': = :J'' J'':
\end{equation}
 由于三个 Sugawara 向量重合，Sugawara 向量的关联函数也彼此重合。后者由 $\mathfrak{sl}(2)$ 量子 Hitchin 哈密顿量作用在配分函数 $\Psi(a)$ 上得到，而 $\Psi(a)$ 被视为 $\M\times \M \times \M$ 上的半密度。 

\begin{yizhu}[译者按：临界 level 与中心元]
仿射 Kac–Moody 代数 $\hat{\mathfrak g}_k$ 的临界 level 为 $k=-h^\vee$，其中 $h^\vee$ 是对偶 Coxeter 数。对 $\mathfrak{sl}(2)$，$h^\vee=2$，故文中 $\widehat{\mathfrak{sl}}(2)_{-2}$ 恰在临界 level；而对 $\mathfrak{sp}(2n)$，$h^\vee=n+1$，所以 §9.1 中 level $-1/2$ 的流远非临界——这正是原文说 $G$ 流的 level “未必临界”、从而可能出现规范反常的原因。临界 level 上 Segal–Sugawara 构造退化：由流双线性型构成的 $:JJ:$ 不再给出 Virasoro 张量，而与所有模式对易、成为中心元；临界 level 真空模具有很大的中心，其结构由 Feigin 与 Frenkel 系统研究（B. Feigin, E. Frenkel, ``Affine Kac--Moody algebras at the critical level and Gelfand--Dikii algebras,'' \emph{Int.\ J.\ Mod.\ Phys.} A7 (1992), Suppl.~1A, 197）。这也说明原文强调“三个 Sugawara 向量重合”为何是有内容的陈述：三者都是中心元，才可以这样逐点比较。
\end{yizhu}

具体地说，辛玻色子 $Z$ 的动能算子是作用在丛 $E \otimes E' \otimes E'' \otimes K_C^{\frac12}$ 上的 $\bar \partial_{a}$ 算子。
配分函数为： 
\begin{equation}\label{sqr}
\Psi(a) = \frac{1}{|\det \bar \partial_a|}.
\end{equation}  
这样，我们就给出了如下交织性质的一个手征代数推导
\begin{equation}
\boxed{H_i \Psi = H'_i \Psi = H''_i \Psi}
\end{equation}
其中 $H_i$ 取遍作用在三份 $\SL(2)$ 丛空间上的 $\mathfrak{sl}(2)$ 量子 Hitchin 哈密顿量。
类似地
\begin{equation}
\boxed{\bar H_i \Psi = \bar H'_i \Psi = \bar H''_i \Psi}
\end{equation}
对共轭的量子 Hitchin 哈密顿量成立，它们作为反全纯微分算子作用在 $\SL(2)$ 丛空间上。
这些关系与 B-模型一侧对三对角界面的预期相符。 
我们的下一个目标是证明 Hecke 算符的类似交织关系。 

为简单起见，我们将使用 $\SO(3)$ 的最小荷 Hecke 算符，尽管辛玻色子耦合到 $\SL(2)$ 丛而非 $\PSL(2)$ 丛。极小 Hecke 修正把 $\SL(2)$ 丛映为被一个 gerbe 扭化的
$\SL(2)$ 丛，反之亦然；因此我们可以一致地描述成对的 Hecke 算符在 $\M(C,\SU(2))$ 上的半密度上的作用。 

一个极小 Hecke 算符会为辛玻色子产生 $Z \to -Z$ 割线的一个端点。这引导我们考虑 Ramond 顶点算子。 

\subsection{Ramond 扇区} \label{sec:Ramond}
辛玻色子手征代数容许 Ramond 模，它们与带有非束缚旋量结构（non-bounding spin structure）的圆周相关联。在这样的模中，$Z^a$ 场的模展开涉及指标 $n$ 为整数的模 $Z^a_n$。
相应的顶点算子引入一条割线，$Z^a$ 跨过它变号。\zy{译者按：圆周上的旋量结构有两种：反周期（bounding）者给出 Neveu--Schwarz 扇区，周期（non-bounding）者给出 Ramond 扇区，故 Ramond 模中模指标为整数。原始文献见 P. Ramond, ``Dual theory for free fermions,'' \emph{Phys.\ Rev.} D3 (1971) 2415；A. Neveu, J. Schwarz, ``Factorizable dual model of pions,'' \emph{Nucl.\ Phys.} B31 (1971) 86。}

零模 $Z^a_0$ 构成一个 Weyl 代数。手征辛玻色子有一族丰富的{\it 最高权} Ramond 模，它们由零模 Weyl 代数的一个模诱导而来。Weyl 模的每个元素都被提升为一个最高权向量/顶点算子，它被正 Fourier 模 $Z^a_n$ 湮灭。负 Fourier 模自由作用，零模则按 Weyl 模中的方式作用。 


Kac–Moody 流的零模 $J_0^{ab}$ 以 $\Sp(2n)$ 生成元 $Z^{(a}_0 Z^{b)}_0$ 的方式作用在一个最高权向量上。所有 Weyl 模都在某种程度上破坏了顶点算子代数（VOA）的 $\Sp(2n)$ 对称性。换言之，不存在带有 $\Sp(2n)$ 不变向量的 Weyl 模。因此，把任何这样的顶点算子插入关联函数，总会在该点减弱 $\Sp(2n)$ 规范不变性。 

产生 Ramond 顶点算子的最简单方式，是考虑 $\PSp(2n)=\Sp(2n)/\Z_2$ 中最小荷的谱流算符。选取一个拉格朗日分裂 $\mathbb{C}^{2n} = V \oplus V^\vee$。取一个奇异规范变换，它在 $V$ 上作用为 $z^{\frac12}$，在 $V^\vee$ 上作用为 $z^{-\frac12}$。所得的谱流算符 $S_V(0)$ 是一个 Ramond 模。$V$ 中 $Z^a$ 的线性组合在趋近 $S_V(0)$ 时按 $z^{\frac12}$ 消失，而 $V^\vee$ 中的线性组合则按 $z^{-\frac12}$ 发散。这意味着 $S_V(0)$ 是一个被 $V$ 中 $Z_0^a$ 的线性组合湮灭的最高权 Ramond 模。特别地，它只依赖于 $V$。 

对任意 $V$，$S_V(0)$ 都可由某个作为参考的 $S_{V_0}(0)$ 经 $\Sp(2n)$ 旋转得到。不失一般性，选取一组基使
\begin{equation}
\omega^{ab} = \delta^{a-b-n}- \delta^{b-a-n},~~a,b=1,\cdots, 2n,\end{equation}
并选取参考顶点算子 $S_{V_0}(z)$ 被 $Z^{n+1}_0, \cdots, Z^{2n}_0$ 湮灭。 
把其余作为产生算子作用的零模记作 $u^a=Z_0^a$（$a\leq n$），并把 $S_{V_0}(z)$ 的相应后代记作 $S_{V_0}[u^a](z)$、$S_{V_0}[u^a u^b](z)$ 等等。湮灭零模则作用为 $Z_0^{a+n}=\partial_{u^a}$。

考虑 Weyl 模中的一个相干态：
\begin{equation}
S_{V_0}[e^{\frac12 B_{ab} u^a u^b}](z)
\end{equation}
它被线性组合 $Z^{a+n}_0 - B_{ac} Z^c_0$ 湮灭。该条件定义一个旋转后的
拉格朗日子空间 $V = B \circ V_0$。于是我们辨认出
\begin{equation}
S_{B \circ V_0}(z) = S_{V_0}[e^{\frac12 B_{ab} u^a u^b}](z)
\end{equation}
确实，我们有 $\partial_{B_{ab}}S_{B \circ V_0}(z) = J_0^{ab}S_{B \circ V_0}(z)$。 

\subsection{非手征 Ramond 模}
接下来，我们可以考虑手征与反手征辛玻色子的联合理论。在 Ramond 扇区中，现在我们有了
手征零模 $Z_0^a$ 与反手征零模 $\bar Z_0^a$ 的作用。由于路径积分围道
把 $\bar Z$ 与 $Z$ 的共轭联系起来，自然要考虑一个使这些
零模互为伴随的 Ramond 态空间。 

此外，零动量扇区中的二维理论本质上是一个以 $\mathbb{C}^{2n}$ 为靶空间的量子力学。
因此自然选取一个极化，并把零模希尔伯特空间取为 $\lmark^2(\mathbb{C}^{n})$。这一答案实际上与极化的选取无关，因为我们可以用广义 Fourier 变换把不同的极化联系起来。\footnote{metaplectic 反常相互抵消，因为我们同时作用于全纯与反全纯变量。} 
完整的 Ramond 希尔伯特空间 ${\mathcal R}$ 按通常方式由 $\lmark^2(\mathbb{C}^{n})$ 诱导而来：当 $k>0$ 时让 $Z^a_k$ 与 $\bar Z^a_k$ 湮灭 $\lmark^2(\mathbb{C}^{n})$ 中的向量，而在 $k<0$ 时让它们自由作用。\zy{译者按：零模对易关系 $[Z^a_0,Z^b_0]=\omega^{ab}$ 给出的正是有限维正则对易关系的 Weyl 代数；选取极化即把 $\mathbb{C}^{2n}$ 分成 $V\oplus V^\vee$，把 $V$ 实现为乘法、$V^\vee$ 实现为微分。不同极化下的表示由 Fourier 变换相联系，这是 Stone--von Neumann 定理的有限维情形。参见 nLab 条目 ``Weyl algebra'' 与 ``Stone--von Neumann theorem''。}

各种不同的最高权 Ramond 顶点算子被实现为 $\mathbb{C}^{n}$ 上的分布。 
选取与 $S_{V_0}(0)$ 的定义中相同的极化。于是参考谱流算符 $S_{V_0}(0)$ 由分布 “$1$” 表示。旋转后的 $S_{B \circ V_0}(z)$ 由高斯型分布
\begin{equation}
e^{\frac12 B_{ab} u^a u^b-\frac12 \bar B_{ab} \bar u^a \bar u^b }
\end{equation}
表示。而分布 $\delta^{(2n)}(u)$ 则表示一个被 $Z^{1}_0, \cdots Z^{n}_0$ 湮灭的 Ramond 顶点算子，
它是荷与 $S_{V_0}(0)$ 相反的谱流算符。

现在我们可以特殊化到 $n=4$，并考虑与 $\Sp(8)$ 的 $\SL(2) \times \SL(2) \times \SL(2)$ 子群相联系的谱流算符。 
我们选取 $\mathbb{C}^{8}$ 的极化，使其在乘积中的第二个与第三个 $\SL(2)$ 群下不变。具体地说，我们取 
$Z_0^{2 \beta \gamma} = u^{\beta \gamma}$ 与 $Z_0^{1 \beta \gamma}=\epsilon^{\beta\beta'}\epsilon^{\gamma\gamma'}\frac{\partial}{\partial u^{\beta'\gamma'}} .$

在这个框架中，三个 $\SL(2)$ 的谱流算符 $\Sigma_{1;\mu, \bar \mu}$、$\Sigma'_{1;\mu', \bar \mu}$ 与 $\Sigma''_{1;\mu'', \bar \mu}$ 都是 $\Sp(8)$ 最小荷谱流算符的特例，即都是最高权 Ramond 顶点算子的特例。它们由分布表示；我们可以对 $\mu$ 平均，从而得到 $\hat \Sigma_{1}$、$\hat \Sigma'_{1}$、$\hat \Sigma''_{1}$ 的表示。计算是直接的：
\begin{itemize}
\item $\Sigma_1$ 由分布 “1” 表示。$\Sigma_{1;\mu, \bar \mu}$ 由高斯型分布
\begin{equation}
e^{\mu  \epsilon_{\beta \beta'}\epsilon_{\gamma \gamma'} u^{2\beta \gamma} u^{2 \beta' \gamma'}- \mathrm{c.c.}}.
\end{equation} 对 $\mu$ 平均，$\hat \Sigma_{1}$ 由分布
\begin{equation}
\delta^{(2)}(\epsilon_{\beta \beta'}\epsilon_{\gamma \gamma'} u^{2\beta \gamma} u^{2 \beta' \gamma'}).
\end{equation}
表示。
\item $\Sigma'_1$ 由分布 
\begin{equation}
\delta^{(2)}(u^{211})\delta^{(2)}(u^{212}).
\end{equation}
 表示。$\Sigma'_{1;\mu', \bar \mu'}$ 由分布 
\begin{equation}
\delta^{(2)}(u^{211}- \mu' u^{221})\delta^{(2)}(u^{212}- \mu' u^{222}).
\end{equation}
 表示。对 $\mu'$ 平均，我们发现 $\hat \Sigma'_{1}$ 由与 $\hat \Sigma_{1}$ 相同的分布表示。
\item $\Sigma''_1$ 由分布 
\begin{equation}
\delta^{(2)}(u^{211})\delta^{(2)}(u^{221}).
\end{equation}
类似地，$\Sigma''_{1;\mu'', \bar \mu''}$ 由分布 
\begin{equation}
\delta^{(2)}(u^{211}- \mu'' u^{212})\delta^{(2)}(u^{221}- \mu'' u^{222}).
\end{equation}
 对 $\mu''$ 平均，我们发现 $\hat \Sigma''_{1}$ 由与 $\hat \Sigma_{1}$ 相同的分布表示。
\end{itemize}
我们由此断定
\begin{equation}
\hat \Sigma_{1} = \hat \Sigma'_{1} =\hat \Sigma''_{1}
\end{equation}
在三基本辛玻色子的理论中成立。这些关系所起的作用与 (\ref{eq:sugathree}) 类似：把它们插入关联函数，就证明了配分函数 $\Psi(a)$ 交织三个 $\SL(2)$ 群的极小 Hecke 算符的作用：
\begin{equation}
\boxed{H_z \Psi =  H'_z \Psi = H''_z \Psi}
\end{equation}

\subsection{一个奇妙模}\label{marvel}
值得详细说明由分布
\begin{equation}
M=\delta^{(2)}(\epsilon_{\beta \beta'}\epsilon_{\gamma \gamma'} u^{2\beta \gamma} u^{2 \beta' \gamma'}).
\end{equation}
生成的 Weyl 模的性质。我们从 $M$ 表示平均后的谱流算符这一事实学到了一些令人惊讶的性质：
$M$ 在 $\SL(2) \times \SL(2) \times \SL(2)$ 下不变，且该模以完全对称的方式对待三个 $\SL(2)$ 群。后一性质在分析中有些隐蔽，因此我们可以在此详细说明：我们可以通过 Fourier 变换改变极化，转到用 $u^{\alpha 2 \gamma}$ 或 $u^{\alpha \beta 2}$ 的函数来表示该 Weyl 模。$M$ 的 Fourier 变换分别产生分布 $\delta^{(2)}(\epsilon_{\beta \beta'}\epsilon_{\gamma \gamma'} u^{\beta 2 \gamma} u^{\beta' 2\gamma'})$ 与 $\delta^{(2)}(\epsilon_{\beta \beta'}\epsilon_{\gamma \gamma'} u^{\beta \gamma 2} u^{\beta' \gamma' 2})$。 

用 Weyl 代数作用于 $M$，我们生成一个带有 $\SL(2) \times \SL(2) \times \SL(2)$ 的显式作用的奇妙模（marvelous module）。
作为向量空间，该模分解为
\begin{equation}\label{marmod}
\oplus_{d=1}^\infty V_d\otimes V_d \otimes V_d
\end{equation}
其中 $V_d$ 是 $\SL(2)$ 的 $d$ 维不可约表示。Weyl 生成元作用为两项之和：一项使 $d$ 增加 $1$，另一项使它减少 $1$。\zy{译者按（无外部依据）：原文未定义 “奇妙” 一词，它只是形容该模的三条性质：由单一分布 $M$ 经 Weyl 代数生成；对三个 $\SL(2)$ 完全对称；作为向量空间恰为 $\oplus_d V_d\otimes V_d\otimes V_d$。}

\subsection{若干推广}
$\SU(2) \times \SU(2) \times \SU(2)$ 规范群可以看作 $\Sp(2n) \times \mathrm{Spin}(2n+2)$ 的一个特例。 
三基本超多重态是 $\Sp(2n) \times \mathrm{Spin}(2n+2)$ 的双基本超多重态的一个特例。
双基本超多重态构造出 $\Sp(2n)$ 与 $\mathrm{Spin}(2n+2)$ 四维规范场论之间的一个 “NS5” 界面。 

双基本辛玻色子理论中相应 Kac–Moody 流的 level 是临界的，因此该 NS5 界面
与变形 Neumann 边界形成非反常的角，我们的构造适用。 

该 NS5 界面 $S$-对偶于 $\mathrm{\Spin(2n+1)}$ 与 $\mathrm{Spin}(2n+2)$ 四维规范场论之间的一个 “D5” 界面，
在后者处 $\mathrm{Spin}(2n+2)$ 规范群被约化为 $\mathrm{\Spin(2n+1)}$。D5 界面将在 B-模型中下降为支撑在嵌入 $\mathrm{Spin}(2n+2)$ 平坦联络空间中的 $\mathrm{Spin}(2n+1)$ 平坦联络空间上的 BBB 界面。这给出关于量子 Hitchin 哈密顿量与 't Hooft 算符在 NS5 界面上作用的简单预言。 
本节的计算应当可以推广以验证这些预言。 

双基本超多重态也可以用来构造 $\U(n)$ 四维规范场论之间的一个 “NS5” 界面，该界面的对偶相对简单，且应当适合二维手征代数分析。$\SU(n)$ 子群的 Kac–Moody 流是临界的，而 $\U(1)$ 反常可以由单个复手征费米子抵消，后者在对角 $\U(1)$ 规范对称性下荷为 $(1,-1)$。 
同样，本节的计算应当可以推广以验证这些预言。例如，最小谱流算符
由分布 $\delta^{(2)}(\det u)$ 表示。 
%Up the contour integral/conformal block subtleties, these relations would already be sufficient to 
%discover remarkable identities for correlation functions, such as
%\begin{equation}
%\oint dt_0 \, \oint dt_\infty \, \langle S_{t_0}(z_0)S_{t_\infty}(z_\infty) \cdots \rangle =  \oint dt'_0 \, \oint dt'_\infty \, \langle S'_{t'_0}(z_0)S'_{t'_\infty}(z_\infty) \cdots \rangle
%\end{equation}
%which by singular gauge transformations can be rewritten as 
%\begin{equation}
%\oint dt_0 \, \oint dt_\infty \, g[t_0,z_0;t_\infty,z_\infty](z) \circ \langle \cdots \rangle =  \oint dt'_0 \, \oint dt'_\infty \, g'[t'_0,z_0;t'_\infty,z_\infty](z) \circ \langle \cdots \rangle
%\end{equation}
%where on the left hand side we act on the vertex operators in the correlation function by 
%the singular gauge transformation $g[t_0,z_0;t_\infty,z_\infty](z)$ in the first $\SL(2)$ factor, and 
%on the right hand side we act on the vertex operators in the correlation function by 
%the singular gauge transformation $g'[t'_0,z'_0;t'_\infty,z'_\infty](z)$ in the second $\SL(2)$ factor.
%
%The gauge transformation $g[t_0,z_0;t_\infty,z_\infty](z)$ acts as $\sqrt{\frac{z-z_0}{z-z_\infty}}$ on vectors of the form $Z_0^{1\beta \gamma}  - t_0 Z_0^{2\beta \gamma}$ and 
%$\sqrt{\frac{z-z_\infty}{z-z_0}}$ on vectors of the form $Z_0^{1\beta \gamma}  - t_\infty Z_0^{2\beta \gamma}$. 
%We can solve that as
%\begin{equation}
%g[t_0,z_0;t_\infty,z_\infty](z) =\frac{1}{\sqrt{(z-z_0)(z-z_\infty)}} \left[z -z_0 + \frac{z_0-z_\infty}{t_0-t_\infty}\begin{pmatrix} t_0 & - t_0 t_\infty \cr 1 & - t_\infty \end{pmatrix}  \right]
%\end{equation}
%In genus $0$, $g[t_0,z_0;t_\infty,z_\infty](z) \circ \langle \cdots \rangle$ represents the correlation function with $x_i \to g[t_0,z_0;t_\infty,z_\infty](z_i) x_i$. 
% 
%We now sharpen the relationship to a relationship between single-valued, non-chiral correlation functions. 
%

\noindent{\it 致谢}   
Perimeter Institute 的研究由加拿大政府通过 Industry Canada、并由安大略省通过 Ministry of Research $\&$ Innovation 资助。
EW 的研究部分由 NSF Grant PHY-1911298 资助。我们感谢 D. Baraglia、P. Etingof、E. Frenkel、D. Kazhdan 与 L. Schaposnik 的讨论。
我们也感谢 P. Etingof 仔细阅读手稿并提出若干澄清。
